\documentclass[10pt,a4paper]{amsart}
\usepackage[margin=19mm]{geometry}
\usepackage{amsmath,amssymb,mathtools}
\usepackage[scr=rsfs,cal=euler]{mathalfa}
\usepackage{xfrac}
\usepackage[unicode,hidelinks,bookmarks=false]{hyperref}
\newcommand{\ie}{i.e.\ }
\newcommand{\cf}{cf.\ }
\newcommand{\resp}{respectively\ }
\newcommand{\Subsection}{\subsection}
\providecommand{\nobreakdash}{\nobreak\hbox{-}}
\setlength{\emergencystretch}{2em}
\multlinegap0pt
\usepackage{mathtools}
\usepackage{float}
\usepackage{amssymb}
\usepackage{tikz}
\newtheorem{lemma}{Lemma}
\def\E{\boldsymbol{E}}
\newcommand{\GL}{L}
\newcommand{\Om}{\Omega}
\newcommand{\bfx}{\boldsymbol{x}}
\newcommand{\curlvec}{\boldsymbol{\mrm{curl}}\,}
\newcommand{\dsp}{\displaystyle}
\newcommand{\mrm}[1]{\mathrm{#1}}
\title{Trapped modes in electromagnetic waveguides}
\author{Anne-Sophie Bonnet-Ben Dhia, Lucas Chesnel, Sonia Fliss}
\date{Source: arXiv:2512.17763v2 (CC BY 4.0)}
\newcommand{\mH}{\mrm{H}}
\begin{document}
\maketitle
\par\medskip\noindent\textbf{Current collection classification: H1.} Retained author material; this item does not qualify for the advanced H2/H3 collection. See the canonical metadata for the source-specific rationale.\par\medskip


{\sloppy \noindent\emph{Source and licence.} Trapped modes in electromagnetic waveguides. Version arXiv:2512.17763v2 (CC BY 4.0), distributed by arXiv under the Creative Commons Attribution 4.0 International licence, \url{http://creativecommons.org/licenses/by/4.0/}, as recorded in the arXiv licence metadata for this exact version and shown on the abstract page \url{https://arxiv.org/abs/2512.17763v2.} Full licence notice: \path{licenses/arxiv-2512.17763-CC-BY-4.0.txt}.\par\smallskip\par}

\noindent\textit{Editorial scope.} Counted unit: the article's lemma EstimAbove (source0.tex lines 1151--1156). The article's complete original proof of it is delivered with it (source0.tex lines 1157--1212, one \texttt{proof} environment of 56 lines). Additional author definitions and conventions recorded in the correction metadata are retained below as uncounted context. Source caveats and pending qualification issues are recorded separately. Author statements and proofs remain unchanged; the fixed complete author source remains bundled.\par\medskip
\par\medskip\noindent\textit{Retained author context (same fixed source; added for the selected problem).}\par
\noindent Define
\[
\begin{array}{l}
\Pi_x\coloneqq[1;+\infty)\times(0;1)\times(0;1)\\[2pt]
\Pi_y\coloneqq(0;1)\times[1;+\infty)\times(0;1)\\[2pt]
\Pi_z\coloneqq(0;1)\times(0;1)\times[1;+\infty)\\[2pt]
\mathcal{C}\coloneqq(0;1)^3
\end{array}
\]
and consider the tripode waveguide $\Om$ such that
\begin{equation}\label{DefTripode}
\Om\coloneqq\Pi_x\cup\Pi_y\cup\Pi_z\cup\mathcal{C}.
\end{equation}

\[
\GL\coloneqq\{(x,y)\in(0;+\infty)^2\,|\,x<1\mbox{ or }y<1\}=(0;1)\times(0;+\infty)\cup(0;+\infty)\times(0;1)
\]
admits an eigenvalue $\lambda_\bullet$ in its discrete spectrum. Consider $\varphi\in\mH^1_0(\GL)$ a corresponding eigenfunction, \textit{i.e.} a function such that 
\begin{equation}\label{Spectral2D}
-\Delta\varphi=\lambda_\bullet\varphi\quad\mbox{ in }\GL
\end{equation}

With the Krein-Rutman theorem, one proves that $\varphi$ can be chosen such that $\varphi>0$ in $\GL$. Moreover, $\varphi$ satisfies $\varphi(x,y)=\varphi(y,x)$ for all $(x,y)\in \GL$.

Define $\E_p$ such that 
\[
\E_p\coloneqq\left(\begin{array}{c}
E_x\\
E_y\\
E_z
\end{array}\right)
\]
with
\[
E_x(\bfx)=\begin{array}{|ll}
\varphi(y,z)& \mbox{ for }x<1\\
0& \mbox{ for }x>1\\
\end{array},\quad E_y(\bfx)=\begin{array}{|ll}
\varphi(x,z)& \mbox{ for }y<1\\
0& \mbox{ for }y>1\\
\end{array},\quad E_z(\bfx)=\begin{array}{|ll}
\varphi(x,y)& \mbox{ for }z<1\\
0& \mbox{ for }z>1.
\end{array}
\]

Therefore we introduce $\psi\in\mH^1_0(\Om)$ such that
\begin{equation}\label{PbDefPsiBis}
\int_\Om\nabla\psi\cdot\nabla\psi'\,d\bfx=\int_\Om\E_{p}\cdot\nabla\psi'\,d\bfx,\qquad \forall\psi'\in\mH^1_0(\Om),
\end{equation}

and we set $\tilde{\E}_{p}\coloneqq\E_{p}-\nabla\psi$.

\begin{lemma}
There holds
\begin{equation}\label{EstimAbove}
\dsp\int_{\Om}|\curlvec \tilde{\E}_p|^2\,d\bfx=3\lambda_\bullet\int_\GL\varphi^2\,dxdy-6\int_0^1\varphi(1,y)^2\,dy.
\end{equation}
\end{lemma}

\begin{proof}
In $\Om$, we have $\curlvec \tilde{\E}_p=\curlvec \E_p$. Moreover, in $\Pi_x$ we find
\[
\curlvec \E_p=\curlvec \left(\begin{array}{c}
0\\
E_y\\
E_z
\end{array}\right)=\left(\begin{array}{c}
\partial_yE_z-\partial_zE_y\\
-\partial_xE_z\\
\partial_xE_y
\end{array}\right).
\]
Thus we obtain
\[
\begin{array}{rcl}
\dsp\int_{\Pi_x}|\curlvec \E_p|^2\,d\bfx&=&\dsp\int_{\Pi_x}(\partial_xE_y)^2+(\partial_zE_y)^2+(\partial_xE_z)^2+(\partial_yE_z)^2-2\partial_yE_z\partial_zE_y\,d\bfx.
\end{array}
\]
By observing that 
\[
\dsp\int_{\Pi_x}\partial_yE_z\partial_zE_y\,d\bfx=\int_1^{+\infty}\int_0^1 \partial_yE_z(x,y)\,dy\int_0^1\partial_zE_y(x,z)\,dzdx=0,
\]
we get 
\[
\dsp\int_{\Pi_x}|\curlvec \E_p|^2\,d\bfx=2\int_{\GL_+}|
\nabla\varphi|^2\,dxdy
\]
where $\GL_+\coloneqq(1;+\infty)\times(0;1)$. Similarly, we obtain
\[
\dsp\int_{\Pi_y}|\curlvec \E_p|^2\,d\bfx=\dsp\int_{\Pi_z}|\curlvec \E_p|^2\,d\bfx=2\int_{\GL_+}|
\nabla\varphi|^2\,dxdy.
\]
On the other hand, in $\mathcal{C}=(0;1)^3$, we have
\[
\curlvec \E_p=\left(\begin{array}{c}
\partial_yE_z-\partial_zE_y\\
\partial_zE_x-\partial_xE_z\\
\partial_xE_y-\partial_yE_x
\end{array}\right).
\]
Therefore, we find
\begin{equation}\label{TermLemma}
\dsp\int_{\mathcal{C}}|\curlvec \E_p|^2\,d\bfx=3\int_{\square}|
\nabla\varphi|^2\,dxdy-2\dsp\int_{\mathcal{C}} \partial_yE_z\partial_zE_y+\partial_xE_z\partial_zE_x+\partial_xE_y\partial_yE_x\,d\bfx
\end{equation}
with $\square=(0;1)^2$. But 
\[
\dsp\int_{\mathcal{C}} \partial_yE_z\partial_zE_y\,d\bfx=\int_0^{1}\int_0^1 \partial_yE_z(x,y)\,dy\int_0^1\partial_zE_y(x,z)\,dzdx=\int_0^{1}E_z(x,1)E_y(x,1)\,dx=\int_0^1\varphi(x,1)^2\,dx.
\]
By exploiting that $\varphi(x,y)=\varphi(y,x)$ in $\GL$, similarly we get
\[
\dsp\int_{\mathcal{C}} \partial_xE_z\partial_zE_x\,d\bfx=\dsp\int_{\mathcal{C}} \partial_xE_y\partial_yE_x\,d\bfx=\int_0^1\varphi(x,1)^2\,dx=\int_0^1\varphi(1,y)^2\,dy.
\]
With (\ref{TermLemma}), we deduce (\ref{EstimAbove}).
\end{proof}

\end{document}
